\section{Short depth blocks and side terms}\label{sec:sides}

Three side factors can be summed at a power $p\in(2/3,1)$, using
their square budgets and the retained normalized $L^2$ bound for
the fourth input.  For exactly two sides, transferring the
complementary filters recovers the full-pair estimate; the remaining
paired near-band terms require an absolute estimate.  We first
prove a short-block comparison with $C_N(q)$, and then reduce the
two varying sides to fixed stacks.  After summing the products of
their sizes, the transfer coefficient $1/P$ leaves a gain $P^{-1/2}$,
up to small overhead and the factor $(K+2)^{C(q-2)}$ from the
short-block comparison.  Throughout the side argument,
Equation~\eqref{eq:src-17} is an explicit hypothesis for the
blockwise constant linear sides; Section~\ref{sec:linearization}
verifies it uniformly on fresh normalized roots.

\subsection{A mixed-norm maximal estimate}

We record the precise maximal estimate needed below.  All the maximal
operators in this subsection act in the real variable $y$ alone.

\begin{lemma}\label{lem:sides-mixed-maximal}
Let $(X,\mu)$ and $(W,\nu)$ be finite measure spaces.  If
$1<r,s,t<\infty$ and $v\in A_r(\R)$, then
\[
 \|M_y g\|_{L_y^r(v;L_\xi^s(X;L_w^t(W)))}
 \le C_{r,s,t}([v]_{A_r})
       \|g\|_{L_y^r(v;L_\xi^s(X;L_w^t(W)))}.
\]
The constant is independent of the numbers of atoms in the measure
spaces.  In particular, for finite exponents $p,b,h>q\geq1$,
\[
 \|M_q f\|_{L_y^pL_\xi^bL_w^h}
 \le C_{p,b,h,q}\|f\|_{L_y^pL_\xi^bL_w^h}.
\]
\end{lemma}

\begin{proof}
Here $v\in A_a$ means that $v>0$ almost everywhere and
\[
 [v]_{A_a}
 =\sup_J\left(\fint_Jv\right)
            \left(\fint_Jv^{-1/(a-1)}\right)^{a-1}<\infty.
\]
We use the scalar weighted maximal theorem and the following
family-of-pairs form of Rubio de Francia extrapolation: if a family
$\mathcal F$ of nonnegative pairs $(U,V)$ satisfies, for some
$a_0\in(1,\infty)$ and every $v\in A_{a_0}$,
\[
 \int U^{a_0}v
 \le C_{a_0}([v]_{A_{a_0}})\int V^{a_0}v
 \qquad ((U,V)\in\mathcal F),
\]
uniformly in the pair, then the corresponding estimate holds at every
$a\in(1,\infty)$ and every $v\in A_a$.  No operator relation between
$U$ and $V$ is assumed; see
\cite[Theorem~3.9]{CUMP2011}.

Initially let $g\geq0$ be bounded and supported in a bounded interval
in $y$.  All the weighted left-hand sides used in the following
applications are finite: the inner norms are bounded by a constant
times the scalar maximal function of the indicator of that interval.
Thus the version of extrapolation with finite left-hand sides suffices.
The scalar weighted $L^t$ maximal inequality, integrated in $w$, gives
the weighted $L^t$ hypothesis for the family of pairs
\[
 U_\xi(y)=\|M_y g(y,\xi,\cdot)\|_{L_w^t},
 \qquad
 V_\xi(y)=\|g(y,\xi,\cdot)\|_{L_w^t}.
\]
Extrapolation gives their weighted $L^a$ inequality for every $a>1$.
Choose $a=s$ and integrate this inequality in $\xi$.  The resulting
pairs
\[
 U(y)=\|U_\xi(y)\|_{L_\xi^s},\qquad
 V(y)=\|V_\xi(y)\|_{L_\xi^s}
\]
satisfy the weighted $L^s$ hypothesis for every $A_s$ weight.  A second
application of extrapolation, now with target exponent $r$, proves the
first assertion.  This argument never interchanges the three norms and
requires no ordering among $r,s,t$.

For general nonnegative $g$, apply the result to
$\min(g,n)\1_{[-n,n]}(y)$ and use monotone convergence, including
$M_y g_n\uparrow M_y g$.  Absolute values treat general $g$.
Finally apply the first assertion to $g=|f|^q$, with
\[
 r=p/q,\qquad s=b/q,\qquad t=h/q,
\]
and take the $q$th root.  Each exponent is strictly greater than one,
and the power identity for the iterated norms is exact.
\end{proof}

\subsection{Comparison on a short block}

A finite dyadic tree piece is completed by its terminal intervals,
called leaves; they partition its root.  Only internal intervals are
used in its form sum.  Consequently, at a contributing point in an
internal interval $I$, the leaf containing that point is a strict
descendant of $I$.

\begin{lemma}[Short-block comparison]\label{lem:short-block}
Let $2<q<3$.  Let $\mathcal T$ be a finite dyadic tree piece with root
$R$, whose internal summation intervals belong to at most $K$
consecutive depths and to the original range of at most $N$ depths.
Suppose its leaves occur in the same depth block or one depth below
it.  In slot $j\in\{0,1,2,3\}$ let $\mathbf z_j$ be a deterministic
Hilbert-space-valued function supported on $R$, with
\[
 \|\mathbf z_j\|_{2,L}\le s_j
 \qquad\text{for every leaf }L.
\]
Let $z_j^\sigma$ be random scalar inputs such that, for every fixed
finite $a\geq1$ used below,
\[
 \bigl(\E_\sigma|z_j^\sigma(y)|^a\bigr)^{1/a}
 \le C_a\|\mathbf z_j(y)\|.
\]
The constants $C_a$ are independent of the tree and $K$; products of at
most two independent Rademacher randomizations of a vector stack have
this property.  For arbitrary measurable kernel signs
$|\epsilon_I(\sigma)|\le1$, and any subcollection of the internal
intervals,
\[
 \E_\sigma\left|
   \sum_I |I|\epsilon_I(\sigma)
             H_I(z_0^\sigma,z_1^\sigma,z_2^\sigma,z_3^\sigma)
             \right|
 \le (1+C_N(q))C_q
       (1+\log(K+2))^C(K+2)^{C(q-2)}
       |R|\prod_{j=0}^3s_j .
 \tag{16}\label{eq:src-16}
\]
Here $C$ is absolute.  The constant $C_q$ may depend on the fixed
moment bounds, as well as the fixed kernel bounds.
\end{lemma}

\begin{proof}
By homogeneity it suffices to take $s_j=1$.  A zero size gives a zero
input almost everywhere and hence a zero form.  To estimate the expected
absolute value, insert a common $\sigma$-dependent scalar of modulus
one into the kernel signs.  It is therefore enough to estimate the
absolute value of the expected signed sum.

\paragraph{Rectangles of leaf depths.}
For each slot write $k_j(y_j)$ for the depth of its leaf.  There are at
most $K+1$ possible depths.  Pad this list to an interval of consecutive
integers of dyadic length at most $4(K+1)$ and use its dyadic
subdivision.  Partition all four-tuples of leaf depths as follows.

If the minimum $t$ occurs in at least two slots, specify its set of
attaining slots.  Those slots have depth $t$, and all the others have
depth greater than $t$.  Designate two attaining slots $a,b$ in a fixed
order.  If the minimum is unique, designate its slot $a$.  The minimum
and the second minimum have a unique smallest dyadic index interval
whose lower and upper halves, denoted $L,U$, contain them respectively.
Specify the nonempty set of slots other than $a$ whose depth lies in
$U$; all remaining slots have depth greater than every index in $U$.
Choose a designated slot $b$ in that nonempty set and put $t=\min U$.
This gives a disjoint partition into rectangles of leaf masks.

There are $O(1+\log(K+2))$ length levels and a bounded number of slot
assignments at each level.  At a fixed length level and assignment,
write $E_{j,\iota}$ for the union of the leaves selected in slot $j$
at position $\iota$.  The sets $E_{a,\iota}$ are disjoint as $\iota$
varies, and so are $E_{b,\iota}$.  There are $O(K+1)$ positions.
The nonunique-minimum cases are treated as an additional level with
singleton positions.

Fix one rectangle and suppress $\iota$.  Every slot has leaf depth at
least $t$, except possibly slot $a$ in the unique-minimum case.  Call
the former slots regular and the latter slot heavy.  We may restrict
the summation depths to $s<t$.  Indeed, if $s\ge t$, one of the selected
leaves has depth at most $s$ and cannot meet the interior of the
internal summation interval at depth $s$.  This observation also
excludes such leaves after the small coordinate enlargements below,
by the fixed kernel support margin.

\paragraph{A slow lattice and a fast torus.}
Write $r_s$ for the interval length at depth $s$, and choose
$\delta=\beta r_t$, where $\beta>0$ is a sufficiently small fixed
dyadic number depending only on the support margin.  Write
\[
 y_0=\delta(n_0+u),\qquad y_1=\delta(n_1+v),
 \qquad n_0,n_1\in\Z,\quad u,v\in[0,1).
\]
On a progression this gives
\[
 n_j=(1-j)n_0+jn_1,
 \qquad w_j=(1-j)u+jv,
 \qquad y_j=\delta(n_j+w_j).
\]
Take a fixed torus $\T_0=\R/(T_0\Z)$ of sufficiently large length,
with normalized Haar measure.  In slots $0,1$ define
$F_j^\sigma(n,w)=z_j^\sigma(\delta(n+w))\1_{E_j}(\delta(n+w))$
on $0\le w<1$ and set it to zero elsewhere on this torus.  In slots
$2,3$ use the same formula on a fixed bounded lift interval containing
all the required $w_j$, and set it to zero outside that interval.
For example, $T_0=16$ and lift intervals contained in $[-4,4]$ suffice.
With fixed normalization factors, the fast Haar progression integral
then reproduces the original integral: the first two factors restrict
$u,v$ to $[0,1)$, and the other lifts do not wrap around the torus.

Replace the kernel at $\delta(n_j+w_j)$ by its value at the slow
coordinates
\[
 (x,t')=(\delta n_0,\delta(n_1-n_0)).
\]
We use $t'$ here for the real progression parameter, to distinguish it
from the depth index $t$.  The difference of the kernels is
$O(\delta/r_s)$, with the same fixed support enlargement.  The
corresponding normalized summand has the factor $(\delta/r_s)^2$
times a lattice sum and a fast progression integral.

Here and below the lattice form bounds follow from the same
pair-coordinate argument as Lemma~\ref{lem:local-absolute}.  For two
different progression slots the integer pair map is injective onto
its image, and the pair map on the connected torus is Haar-surjective.
Thus the two-$L^1$, two-$L^\infty$ estimates hold on the product of
the normalized local counting measure and the fast torus.  Interpolation
gives the local absolute estimate whenever the four reciprocal
exponents have sum at most two.

For clarity, all kernel-sampling errors can be summed as follows.
Let $M_{j,I}$ be the squared mass of the deterministic stack on the
sampling windows relevant to $I$, including their bounded overlap
multiplicity.  The support margin and the choice of $\beta$ place
these windows inside $I$.  Hence
\[
 M_{j,I}\le C|I|,\qquad
 \sum_{I:\,\operatorname{depth}(I)=s} M_{j,I}\le C|E_j|.
\]
The absolute $L^2$ progression estimate, followed by H\"older in
$\sigma$, bounds the sum of the errors at depth $s$ by
\[
 C\frac{\delta}{r_s}
 \left(\sum_{\operatorname{depth}(I)=s}M_{a,I}\right)^{1/2}
 \left(\sum_{\operatorname{depth}(I)=s}M_{b,I}\right)^{1/2}
 \le C\frac{\delta}{r_s}|E_a|^{1/2}|E_b|^{1/2}.
\]
We used the normalized mass bounds in the other two slots.  Since
$\sum_{s<t}\delta/r_s\le C$, these errors cost at most
$C|E_a|^{1/2}|E_b|^{1/2}$ for the rectangle.  Summing positions by
the two designated disjointness properties costs $C|R|$ at each
length level.

\paragraph{Quadratic randomization on the fast torus.}
Put $c=(-1,3,-3,1)$.  On Fourier coefficients in the fast variable
define the unitary maps
\[
 \widehat{U_{j,\gamma}F}(n,m)
   =e(\gamma m^2/c_j)\widehat F(n,m),
 \qquad \gamma\in\R/(3\Z).
\]
The fast progression integral is unchanged when all four factors are
replaced by $U_{j,\gamma}F_j$.  Indeed its Fourier expansion is
supported on
\[
 \sum_jm_j=\sum_jjm_j=0.
\]
The nullspace of these two equations is spanned by $(c_j)_j$ and
$(jc_j)_j$, so $m_j=c_j(h+jh')$ for real $h,h'$.  Since
$\sum_j j^\ell c_j=0$ for $\ell=0,1,2$, we have
\[
 \sum_jm_j^2/c_j=0.
\]
Moreover, for each single slot and each fixed $n,\sigma$,
\[
 \E_\gamma\int_{\T_0}|U_{j,\gamma}F_j^\sigma(n,w)|^4\,dw
 \le 2\|F_j^\sigma(n,\cdot)\|_2^4.
\]
To see this, integration in $w$ equates the sums of two Fourier
indices, and integration over the common period $3$ equates their
sums of squares.  The two pairs are consequently equal as multisets.
The two possible pairings give the displayed bound.  The identities
are first proved for Fourier polynomials.  The pairwise $L^2$
progression estimate and the fourth-moment bound, also applied to
differences, justify passage to $L^2$ limits.

Let $e_j(n)$ be the fast $L^2$ norm of the deterministic stack on the
corresponding sampling window.  Minkowski and the pointwise sign
moment hypothesis give, for every needed finite $a$,
\[
 \bigl(\E_\sigma\|F_j^\sigma(n,\cdot)\|_2^a\bigr)^{1/a}
 \le C_a e_j(n).
\]
The windows overlap a bounded number of times, so
\[
 \delta\sum_n e_j(n)^2\le C|E_j|.
\]
If slot $j$ is regular, then $e_j(n)\le C$.  In fact, all selected
leaves have length at most $r_t$.  Disjoint leaves meeting a window of
length $O(\delta)$ have total length at most $O(\delta)+2r_t$: all but
at most two are contained in that window.  Summing their complete
leaf masses and dividing by $\delta$ proves the bound, with a constant
depending only on the fixed $\beta$.  On the lattice indices relevant
to any summation interval $I$, the local version of the preceding mass
estimate is
\[
 \frac{\delta}{|I|}\sum_{n\ \mathrm{relevant\ to}\ I}e_j(n)^2
 \le C
\]
in all four slots, including the heavy slot.  All the required windows
are inside $I$, whose leaves have normalized squared mass at most one.

\paragraph{Clipping the possible heavy slot.}
Write $G_j^{\gamma,\sigma}=U_{j,\gamma}F_j^\sigma$.
If a heavy slot is present, clip $G_a$ radially at
$A=(K+2)^{10}$; write the clipped function as $G_a^A$.
For the normalized local counting measure in an interval $I$ and
Haar measure in $w$, unitarity gives, at every $\gamma$,
\[
 \|G_a-G_a^A\|_1
 \le A^{-1}\|G_a\|_2^2
 =A^{-1}\|F_a^\sigma\|_2^2.
\]
Its $L^4$ norm in the random variables is $O(A^{-1})$, by the local
mass estimate and the eighth sign moment.  For each regular slot,
the local $L^3$ norm of $G_j$ has bounded fourth moment in
$(\gamma,\sigma)$.  Indeed the local measure is bounded, and the
fourth-moment identity bounds its fourth power by a constant times
\[
 \frac{\delta}{|I|}
       \sum_{n\ \mathrm{relevant\ to}\ I}e_j(n)^4\le C.
\]
The absolute progression estimate with exponents $(1,3,3,3)$ and
H\"older of exponent four in the shared random variables therefore
bound the expected tail error by $C/A$ per normalized cell.  Even
the crude sum over $O(K+1)$ positions and $O(K+1)$ depths costs only
$C(K+2)^2/A$ times $|R|$ per length level.  This is bounded uniformly
in $K$.  Henceforth $G_a$ denotes the clipped function when there is
a heavy slot.

\paragraph{Return to the continuous form by bursts.}
For fixed $\gamma,\sigma,u,v$, with $u,v$ now ranging over the whole
fast torus, set
\[
 g_j(y)=\sum_n G_j^{\gamma,\sigma}
                (n,(1-j)u+jv)\,
                \1_{\{|y-\delta n|<c\delta\}}(y),
\]
where $c>0$ is a sufficiently small fixed constant.  The torus
argument in this formula is taken modulo $T_0$.  If a real progression
meets four such bursts, then the integer centers satisfy exactly
$n_j=(1-j)n_0+jn_1$.  For example, the discrepancies for $j=2,3$ are
integers of absolute value less than one when $c<1/20$.
For a fixed slow pair, the admissible offset polygon
\[
 \{(a_0,a_1): |(1-j)a_0+ja_1|<c, 0\le j\le3\}
\]
has a fixed positive area.  Thus, after replacing the continuous
kernel by its value at the center, its integral on these bursts is
that fixed area times $\delta^2$ times the desired sampled product.

This second replacement has the same summed error bound as the first
one.  On the relevant indices, all the original fast sampling windows
still lie inside $I$.  Unitarity and clipping bound the transformed
local $L^2$ sizes by those of $F_j^\sigma$.  The lattice-times-torus
pairwise estimate, then H\"older in the random variables, therefore
uses the same masses $M_{j,I}$ and the same geometric factor
$\delta/r_s$.  Although Fourier evolution spreads the fast support
on the torus, it does not create new slow indices, so this support
argument remains valid.

We may now apply the definition of $C_N(q)$ to the burst inputs with
the original admissible kernels and original summation intervals.
Only the comparison error used sampled kernels; no marginal
cancellation hypothesis has been changed.  It remains to bound
\[
 C_N(q)\int_\R\E_{\gamma,\sigma}
       \int_{\T_0^2}\prod_{j=0}^3 M_qg_j(y)\,du\,dv\,dy.
\]

\paragraph{The exponents and the dependence on $K$.}
Put
\[
 p_a=p_b=p_*=2+2(q-2),\qquad
 p_j=\frac{2p_*}{p_*-2}\quad(j\notin\{a,b\}).
\]
These four exponents are finite, strictly greater than $q$, and have
reciprocals summing to one.  Use fast exponents $h_a=p_*$ for a heavy
slot, and $h_j=4$ for every regular slot.  The sum of their reciprocals
is at most two.  The fast progression H\"older estimate, followed by
H\"older of exponent four in $(\gamma,\sigma)$ and spatial H\"older
with the exponents $p_j$, bounds the last display by
\[
 C_q C_N(q)\prod_{j=0}^3
 \|M_qg_j\|_{L_y^{p_j}L_{\gamma,\sigma}^4L_w^{h_j}}.
\]
All three exponents in each slot exceed $q$.  Apply
Lemma~\ref{lem:sides-mixed-maximal} in exactly this norm order.

For a regular slot the fourth-moment estimate gives
\[
 \|g_j\|_{L_y^{p_j}L_{\gamma,\sigma}^4L_w^4}
 \le C\left(\delta\sum_ne_j(n)^{p_j}\right)^{1/p_j}
 \le C|E_j|^{1/p_j}.
\]
For the heavy slot, clipping and unitarity give pointwise in
$(\gamma,\sigma)$
\[
 \|G_a(n,\cdot)\|_{p_*}
 \le A^{1-2/p_*}\|F_a^\sigma(n,\cdot)\|_2^{2/p_*}.
\]
Taking the fourth random moment and then the spatial $p_*$ norm yields
\[
 \|g_a\|_{L_y^{p_*}L_{\gamma,\sigma}^4L_w^{p_*}}
 \le C_q A^{1-2/p_*}|E_a|^{1/p_*}.
\]
No dimension or number of signs occurs in these bounds.

At a fixed length level and assignment there are $J\le C(K+1)$
positions.  Designated disjointness and H\"older give
\[
 \sum_{\iota=1}^J
    |E_{a,\iota}|^{1/p_*}|E_{b,\iota}|^{1/p_*}
 \le J^{1-2/p_*}|R|^{2/p_*}.
\]
The other masks have length at most $|R|$, so their factors contribute
$|R|^{1-2/p_*}$.  The resulting power of $K+2$ is at most
\[
 A^{1-2/p_*}(K+2)^{1-2/p_*}
 =(K+2)^{11(1-2/p_*)}
 =(K+2)^{11(q-2)/(q-1)}.
\]
The number of length levels is logarithmic; the two sampling errors
and the clipping errors have already been bounded at that cost.
This proves Equation~\eqref{eq:src-16}, with an absolute exponent
$C$.  All Fourier and measurability approximations are harmless:
there are finitely many slow indices, the fourth-moment estimate
controls differences, clipping is continuous in $L^2$, and the
mixed-norm estimate gives convergence of the final bounds.
\end{proof}

\subsection{The side hypotheses and variation of coarse outputs}

Return to one attempt on a normalized root $R$, one fixed active
pattern in Equation~\eqref{eq:src-11}, and one search leaf in
Equation~\eqref{eq:src-14}.  Write $D_{\mathrm m}$ for the main
catalog budget.  Thus $\log D_{\mathrm m}\le k^{O(1)}$, the number
$\mathcal L$ of colors is $k^{O(1)}$, and each relevant history has
length $m=O(\log\mathcal L)$.  Write
$G_0=\exp(\log^{O(1)}(2k))$; increasing its fixed exponent or taking
a fixed power is allowed.  Since $K\le\operatorname{poly}(P,k)$,
all factors $(C\log(K+2))^{O(m)}$, as well as the compression factors
with budget $D_{\mathrm m}$, are of this form.

Partition the depths into contiguous blocks $[b,b+K)$ and choose a
common pre-depth $u_b\le b$ with $b-u_b\le C K$.  Outputs before the
attempt are set to zero.  In a basic-$B$ slot put
\[
 S_{l,G,s}^{\mathrm{roll}}=Z_{l,G,s}-Z_{l,G,u_b},
 \qquad b\le s<b+K.
\]
We also allow a function $V_{l,G,b}$ which is fixed throughout that
block, may be piecewise in space, and is supported on branches with
needed outputs or downstream use certificates.  The hypothesis on
this linear side is
\[
 \sum_{G,b}\|V_{l,G,b}\|_2^2\le G_0|R|,
 \qquad G_0=\exp(\log^{O(1)}(2k)).
 \tag{17}\label{eq:src-17}
\]
This is assumed uniformly whenever the construction is started on a
normalized root.  On a fresh attempt the linear sides may be
reconstructed from the inherited basic data; no proportional local
energy bound is being inferred by merely restricting an arbitrary
family satisfying the displayed global inequality.
In a basic-$B$ slot, a side state $S_{l,G,s}$ is a sum of its rolling
side and $V_{l,G,b}$.  In a basic-$U$ slot the entire current output
$Z_{l,G,s}$ is a mandatory side.  The basic-$U$ square bound from
Section~\ref{sec:base} is available on every normalized subroot.
We shall prove the estimates uniformly over these finitely many
choices of side component.

The helper projections in Lemma~\ref{lem:transfer} are constructed
after these entire side states have been fixed.  Their individual
mask projections are nested in depth, including their inherited
values between actual uses.  Their histories have length $O(m)$.
Only whole states and those histories are used here; no choice of a
mode, height, or amplitude inside a side state is a helper base.

\begin{lemma}[Variation on certified endpoint cells]
\label{lem:sides-variation}
Let $Z_{G,s}$ be a basic coarse output in
Equation~\eqref{eq:src-14}.  Consider a common collection of disjoint,
or boundedly overlapping, depth intervals $[r,t]$.  At each right
endpoint, restrict to a union of $t$-cells having a downstream actual
use certificate for $G$.  Then
\[
 \sum_{G,[r,t]}\|
  \1_{E_{G,r,t}}(Z_{G,t}-Z_{G,r})\|_2^2
 \le G_0\|\1_RF\|_2^2.
\]
The same conclusion holds for a shared coarse history with no outer
ball, omitting the group sum.  The constant only acquires the given
overlap multiplicity.
\end{lemma}

\begin{proof}
For this estimate expand complements by linearity and couple each
pad choice across all times in an average.  Jensen permits us to
treat a fixed choice.  We have
\[
 Z_{G,s}=L_{G,s}B'_sF,
\]
where $L_{G,s}$ is an individual nested ball projection and $B'_s$
is a shared product of $O(m)$ nested main projection families.  Basic
$B$ inputs are fixed signed sums of at most two such current
histories, which changes only a fixed factor.  Split
\[
 Z_{G,t}-Z_{G,r}
 =(L_{G,t}-L_{G,r})B'_rF
       +L_{G,t}(B'_t-B'_r)F.
\]

For a fixed group, the rows
$\1_{E_{G,r,t}}(L_{G,t}-L_{G,r})$ have a square Bessel bound, since
increments of a nested orthogonal projection family on disjoint
depth intervals are orthogonal.  Their adjoints have ranges in sums
of whole ball ranges at the two endpoints.  A right-endpoint
certificate also certifies the required earlier input cell when its
mask is nonempty.  Splitting at the $r$-cells if necessary therefore
places these adjoint ranges in the usable endpoint spaces of
Lemma~\ref{lem:usable-bessel}.  That lemma makes the rows jointly
Bessel across groups.  They are local to their $r$-cells, so
Lemma~\ref{lem:compression} inserts the shared product $B'_r$ at a
cost $G_0$.

For the second term, the usable-range Bessel estimate at each common
right endpoint first sums the outer ball projections over groups.
Telescope $B'_t-B'_r$.  In each telescoping term, discard the left
contractions.  The changed middle factor is an orthogonal increment
on $[r,t]$, and the factors to its right are evaluated at $r$.
These middle increments form Bessel rows, and
Lemma~\ref{lem:compression} inserts their right factors.  Summing
the $O(m)$ telescoping terms is again absorbed in $G_0$.

An interval coloring treats bounded overlap.  Expansions of the
$O(m)$ complements cost $2^{O(m)}$, and averages are restored by
Jensen, without a factor counting pad choices.  The proof without
an outer ball is the same middle-increment argument.  Intervals
crossing the start of the attempt occur with bounded overlap; the
zero convention and the one-time projection bound treat them.
\end{proof}

The integrated local squared sizes of all the side components obey
\[
 \sum_{s,G}\sum_{I\ \mathrm{used\ at}\ s}
       |I|\|S_{l,G,s}\|_{2,I}^2\le P^C|R|.
\]
For a basic-$U$ side, current group Bessel and the basic square bound
give this even with $G_0$ on the right.  For a rolling side apply
Lemma~\ref{lem:sides-variation} to $[u_b,s]$; these intervals have
$O(K)$ overlap.  A fixed $V_{l,G,b}$ is counted at most $K$ times,
so Equation~\eqref{eq:src-17} gives the assertion for it.
First-crossing stops for the partial path sums consequently give,
outside subroots of arbitrarily small fixed total relative length,
\[
 \sum_{s,G}\|S_{l,G,s}\|_{2,I_s(x)}^2\le P^C
\]
on retained paths.  The fixed exponent can be increased to impose
all polynomially many power-cost caps used below.  In this notation
only actual uses are counted.  The same convention applies to shared
side data, with no group index.

\subsection{Power estimates and frozen complements}

We explain which terms can be estimated before using the stronger
stack argument.  Suppose that a term has three strong size factors
and a fourth bounded factor.  For two mandatory group-indexed sides
write, along a path,
\[
 a_s=\Bigl(\sum_Ga_{s,G}^2\Bigr)^{1/2},\qquad
 b_s=\Bigl(\sum_Gb_{s,G}^2\Bigr)^{1/2}.
\]
The third size is either shared or is bounded by its analogous
group-aggregated size $c_s$.  Cauchy--Schwarz in the group index
bounds the absolute group sum by $P^C a_sb_sc_s$; the factor $P^C$
also covers a retained supremum bound for the fourth input.  For
$2/3<p<1$,
\[
 \sum_s(a_sb_sc_s)^p
 \le
 \Bigl(\sum_sa_s^{3p}\Bigr)^{1/3}
 \Bigl(\sum_sb_s^{3p}\Bigr)^{1/3}
 \Bigl(\sum_sc_s^{3p}\Bigr)^{1/3}
 \le
 \Bigl(\sum_sa_s^2\Bigr)^{p/2}
 \Bigl(\sum_sb_s^2\Bigr)^{p/2}
 \Bigl(\sum_sc_s^2\Bigr)^{p/2}.
\]
The path caps and integration in the path point give a total
$P^C|R|$ at power $p$.  A fourth strong factor may use its retained
supremum.  In particular, this estimate takes the power after the
group sum; it does not apply the cell minimum separately to each
group.

Consider a two-side term with whole current complements as in
Lemma~\ref{lem:transfer}.  On removing the complementary filters,
the main term has inputs
\[
 (\mathcal D_{a,s}S_{a,G,s},S_{b,G,s},B_c,B_d).
\]
Replacing $B_c,B_d$ by the full $Y_c,Y_d$ makes errors with a third
basic side.  The full complementary pair then has norm at most
$\eta$, by the base construction.  Local contractions on the first
side and Cauchy--Schwarz in the two side sizes give total main cost
$\eta P^C|R|$.  The numerical errors of the transfer identity have
an arbitrarily small inverse base-budget factor times those same
two sizes and are negligible.  The only remaining terms are the
coupled near bands of Equation~\eqref{eq:src-15}.

In a block, denote the outgoing complementary band stack by
$(C_{G,s}^q)_q$ and the other complementary input by $D_{G,s}$.
Freeze both at $b$.  Current pad differences are mutually orthogonal
projections, so the current local stack norm of $(C_{G,s}^q)_q$ is
bounded; the current norm of $D_{G,s}$ is bounded as well.  Their
differences from their values at $b$ have squared budgets $P^C|R|$,
by Lemma~\ref{lem:sides-variation}, paying the polynomial number of
band indices and the $O(K)$ overlap.  When the last complementary
band is $1-Q_P$, split it into an identity term and a positive ball
term, as in Lemma~\ref{lem:transfer}.  The positive outer-ball
terms have joint group budgets, and the identity or cumulative
histories have shared budgets.  Thus this argument does not count
the same shared history once for every group.

Expanding a freeze error gives a third strong factor.  If its other
frozen complement is unbounded, write that complement as its bounded
current value plus its strong difference.  The preceding power
estimate still applies.  The same observation allows us to discard
uses at which a frozen complement is large, either at that use or
at an ancestor in its block.  More explicitly, if a current stack
has size at most $C_0$ and its frozen stack has size greater than
$2C_0$, their difference has size greater than $C_0$; hence the
indicator of this event is bounded by $C_0^{-2}$ times the squared
difference size.  Lemma~\ref{lem:sides-variation} applies at all
potential ancestor cells with downstream certificates.  Propagating
such an event down its block costs at most $K$ in the integrated
square budget.  Its square-root indicator is therefore a third
strong factor.  Include these budgets in the power-cost path caps.
On all remaining uses, both frozen complements have bounded local
$L^2$ norms on their entire ancestor hull inside the block, using
the stack norm for the banded complement.

\subsection{Fixed stacks for the two sides}

\begin{lemma}[Side stacks]\label{lem:sides-stacks}
Fix a near offset $q-p$ in Equation~\eqref{eq:src-15}, a transfer
step, and a search path.  The two time-dependent side inputs can be
written, at their uses, as a sum of at most small-overhead many
configurations of fixed vector stacks for each $(G,b)$.  One
position is selected at a use from each stack.  The banded side
stack also includes the band index.  After absorbing the number of
configurations in $G_0$, their energies satisfy
\[
 \sum_{G,b}\|\mathrm{stack}_{a,G,b}\|_2^2
       \le P G_0|R|,
 \qquad
 \sum_{G,b}\|\mathrm{stack}_{\mathrm{untouched},G,b}\|_2^2
       \le G_0|R|.
 \tag{18}\label{eq:src-18}
\]
The functions in these stacks are fixed before subsequent cuts of
the summation set.
\end{lemma}

\begin{proof}
For a sequence $W_s$, $b\le s<b+K$, write
\[
 W_s=W_b+\sum_{[r,t]\in\mathcal P(b,s)}(W_t-W_r),
\]
where $\mathcal P(b,s)$ is the usual disjoint binary decomposition
of the prefix $[b,s]$.  There is at most one interval of each length
in this decomposition.  For a fixed length, its intervals are
disjoint as depth intervals, and the selected endpoint is at most
$s$.  Restrict each resulting fixed function to the union of the
use cells selecting it.  This preserves its value at every required
use and only decreases its norm.  The base position and the choices
of lengths account for logarithmic factors.

For a basic-$U$ input, use each depth $s$ directly as a position.
The current side square bound controls the sum of its energies;
subsequent helper products are local contractions.  On the banded
side, paying at most $P+1$ indices is permitted in the first
inequality of Equation~\eqref{eq:src-18}.

For a fixed linear side $V_{G,b}$, let $T_s$ be the product of its
subsequent averaged projections.  At a fixed band index, the outer
band difference is a difference of two such products.  The
binary-prefix product estimate in the proof of
Lemma~\ref{lem:compression}, now on a sequence of at most $K+1$
times, gives
\[
 \|T_bV_{G,b}\|_2^2+
 \sum_{[r,t]\ \mathrm{at\ a\ fixed\ length}}
       \|(T_t-T_r)V_{G,b}\|_2^2
 \le (C\log(K+2))^{O(m)}\|V_{G,b}\|_2^2.
\]
Individual projections are nested; averages are coupled across
times and restored by Jensen.  Sum the length levels and use
Equation~\eqref{eq:src-17}.  Pay at most $P+1$ on the banded side.

For a rolling input write $W_s=T_s(Z_{G,s}-Z_{G,u_b})$.  The base
term is controlled by
Lemma~\ref{lem:sides-variation} on the intervals $[u_b,b]$, whose
overlap across blocks is bounded.  Its increments split exactly as
\[
 W_t-W_r
 =T_t(Z_{G,t}-Z_{G,r})
       +(T_t-T_r)(Z_{G,r}-Z_{G,u_b}).
\]
The first term is bounded directly by
Lemma~\ref{lem:sides-variation}.  When estimating norms, keep the
entire endpoint cells containing the selected use cells: these
cells still have downstream certificates, and locality of $T_t$
makes this enlargement legitimate.

For the second term, also expand $[u_b,r]$ by binary prefix intervals
in a common window of $O(K)$ depths.  Fix an outer length, an inner
length, and an inner interval $[a,v]$.  For all outer positions
$[r,t]$ using this same inner interval, use the fixed input
\[
 X=\1_{\Omega}(Z_{G,v}-Z_{G,a}),
\]
where $\Omega$ is the union of the needed $r$-cells.  This input is
independent of the outer position.  Since $v\le r$, locality on
$r$-cells shows that it gives the required outputs of $T_t-T_r$.
Its norm is bounded by keeping the union of the $v$-cells containing
those $r$-cells; all of these have downstream use certificates.
For this fixed input, the product-jump estimate bounds the sum over
outer positions by $(C\log(K+2))^{O(m)}\|X\|_2^2$.
Now sum the squared inputs over groups and inner intervals at the
fixed inner length, using
Lemma~\ref{lem:sides-variation}.  The windows for different blocks
overlap a bounded number of times.  Finally sum the two length
choices.  This proves the rolling estimate with only small overhead,
and with the allowed factor $P$ for a banded side.

The number of search leaves, transfer steps, choices of side
component, and length configurations is small overhead.  Individual
pad choices are not additional configurations: they remain inside
the actual averaged operators throughout.  Thus summing all required
configurations preserves the asserted bounds after increasing
$G_0$.
\end{proof}

The reduction to fixed stacks does not require synchronizing their
two position decompositions.  For each fixed configuration choose
independent Rademacher signs $(\varepsilon_i)$ and $(\eta_j)$ for
the two position indices, and a third independent family $(\zeta_p)$
for the paired band indices.  Randomize the banded fixed inputs by
$\varepsilon_i\zeta_p$, the untouched side inputs by $\eta_j$, and
the frozen complementary band stack by $\zeta_p$.  At a use, multiply
the kernel by the selected position signs
$\varepsilon_{i(s)}\eta_{j(s)}$.  Expectation selects exactly the
two required positions and the matched bands.  The remaining
coefficient is $(p-q)/P=O(1/P)$.  Iterated Khintchine gives every
fixed pointwise moment bound required in
Lemma~\ref{lem:short-block}, with the deterministic stack norms.
Initial kernel signs may be arbitrary.

\subsection{Hull maxima and common cuts}

We give the support convention needed to use
Equation~\eqref{eq:src-18} without an additional power of $P$.
For each index $\nu=(G,b,\text{fixed configuration})$, let
$\Omega_\nu$ be the union of its use cells.  Its maximal use cells
are disjoint roots.  The hull consists of the cells between those
roots and their uses.  All fixed stack functions are kept unchanged
when uses are subsequently removed.  A maximal function on a hull
is defined on the entirety of each of its intervals, including
missing branches.

The sum of the lengths of $\Omega_\nu$ is at most $M_0|R|$, where
$M_0\le\exp(k^{O(1)})$.  This follows from the active path
multiplicity, the main group count, and the small-overhead number of
configurations.  Suppose one family of stacks has total energy at
most $E|R|$.  Its squared hull maxima
\[
 m_\nu(x)=\sup_{\substack{I\ \mathrm{in\ its\ hull}\\x\in I}}
                       \fint_I\|\mathrm{stack}_\nu\|^2
\]
satisfy, by the weak $L^1$ inequality for the dyadic maximal operator,
\[
 \sum_\nu|\{m_\nu>T\}|
 \le\min\bigl(M_0|R|,CE|R|/T\bigr).
\]
First stop at individual crossings of $C_1E$, losing at most
$C/C_1$ of the root length.  Integrating the last distributional
bound up to $C_1E$ gives
\[
 \sum_\nu\int_R\min(m_\nu,C_1E)
 \le CE\bigl(1+\log(2+C_1M_0)\bigr)|R|.
\]
Markov's inequality then gives a joint square cap at a further
small fixed loss of root length.  The logarithm, and the constants
needed for the small-overhead many configurations, are absorbed in
$G_0$.  Applying this to the two energies in
Equation~\eqref{eq:src-18} gives caps $PG_0$ and $G_0$ respectively.

These stops are made on partial ancestor suprema, not merely on a
set specified at terminal points.  More precisely, at depth $s$
sum the suprema over the hull ancestors of depth at most $s$ and
stop at the first threshold crossing.  Remove that cell and its
descendants from every original summation set.  Each triggering
cell is removed before its value is included in a retained cap.
Consequently the maxima of the retained hulls satisfy the joint
caps everywhere, even inside a stopped region: only its earlier
retained ancestors still contribute there.  This construction is
dominated by the preceding full-tree distributional bounds and
has the same measure cost.  Combine these cuts with the first-crossing
cuts for the partial power-cost budgets, including the budgets of the
large-frozen-complement flags, by taking their union and using the
original summands at each first cut.  A use flagged for a large frozen
complement is removed only from the critical near-band sum and is
charged to the powered error estimate; the flag itself does not create
a restart root.

\subsection{The near-band estimate and the conditional side bound}

Fix one of the configurations in Lemma~\ref{lem:sides-stacks}.
On each retained hull, make a sparse stopping decomposition for the
two fixed side stacks.  From a piece root $Q$, stop at the first
descendants where either squared stack average exceeds sixteen
times its average on $Q$.  Their total length is at most
$|Q|/8$.  If a root average is zero, the corresponding stack is
zero on $Q$ and there is no contribution there.  The sets $E_Q$
left after removing the next roots have measure at least
$7|Q|/8$ and are disjoint within each fixed hull construction.

Complete every piece by children immediately below its included
cells.  This includes branches without a use and boundaries made
by the common cuts.  A child's squared average is at most twice
its parent's.  Thus both side stacks have bounded normalized
sizes, relative to their root sizes, on every leaf of the piece.
The frozen complementary stack and the other frozen complement
have uniformly bounded sizes on retained hull cells, and hence
on these boundary leaves as well.  The piece uses at most $K$
depths, with leaves at most one step below.  Apply
Lemma~\ref{lem:short-block} to the randomized fixed inputs.

Let $a_Q,b_Q$ be the two side stack sizes at the piece root.  On
every $x\in E_Q$, including a missing or stopped branch, the
retained hull maxima satisfy
\[
 a_Qb_Q\le\sqrt{m_{a,\nu}(x)m_{b,\nu}(x)}.
\]
Indeed $Q$ itself is one of the intervals in each supremum.  Hence
\[
 \sum_{\nu,Q}|Q|a_Qb_Q
 \le C\int_R\sum_\nu
              \sqrt{m_{a,\nu}m_{b,\nu}}
 \le C\int_R
       \Bigl(\sum_\nu m_{a,\nu}\Bigr)^{1/2}
       \Bigl(\sum_\nu m_{b,\nu}\Bigr)^{1/2}
 \le P^{1/2}G_0|R|.
\]
The logarithms in Equation~\eqref{eq:src-16} are small overhead.
After the coefficient $O(1/P)$ and the small-overhead number of
configurations, the total near-band contribution is at most
\[
 (1+C_N(q))C_qP^{-1/2}G_0(K+2)^{C(q-2)}|R|.
 \tag{19}\label{eq:src-19}
\]
This bound holds with arbitrary initial signs at the original cells.
It therefore controls the absolute sum of the cell contributions,
provided the prescribed paired band sums at a fixed offset are
retained inside each contribution.  No absolute summation over
their individual band indices was made.

\begin{proposition}[Side estimate]\label{prop:side-estimate}
Assume the base and localization hypotheses above, the helper
conclusions of Lemma~\ref{lem:transfer}, and the linear-side square
bound \eqref{eq:src-17}.  Fix $p\in(2/3,1)$.  After common
first-crossing cuts of arbitrarily small fixed total relative
length, terms having at least three side factors have total
length-weighted $p$th-power cost at most $P^C|R|$, provided the
fourth factor has retained normalized $L^2$ size at most $P^C$.
This includes a fourth side, a bounded current argument, or a
bounded current argument minus side factors.

Terms having exactly two side factors and whole current coarse
complements admit the following estimates.  Their freeze errors,
large-complement terms, and restoration errors have total
length-weighted $p$th-power cost at most $P^C|R|$.  Their transferred
main terms cost at most $\eta P^C|R|$ absolutely.  Their numerical
errors are negligible after the stated choice of the main
precision.  Their remaining coupled near-band terms obey
Equation~\eqref{eq:src-19} without taking a power less than one.

All statements are uniform on normalized subroots with the same
basic data.  Subsequent common cuts may remove original summands
without changing the fixed functions used in these estimates.
\end{proposition}

\begin{proof}
The three-side power bound was proved above using the path square
caps and $3p>2$.  Lemma~\ref{lem:transfer} gives the decomposition
into transferred main terms, numerical errors, and coupled near
bands.  Restoring the full complementary pair in a main term
introduces a third basic side and leaves the pair-small estimate
$\eta$ on its two arbitrary side inputs.  The frozen-complement
argument accounts for every freeze error and every discarded large
complement as a third strong factor.  Lemma~\ref{lem:sides-stacks},
the common maximal cuts, and the sparse comparison then give
Equation~\eqref{eq:src-19} for the remaining terms.

All cuts were imposed on nonnegative partial ancestor quantities
and removed the original summands at and below their first
violation.  The normalized base hypotheses, the basic square
estimate, and the root-uniform hypothesis \eqref{eq:src-17} are
available anew when the construction is restarted on such subroots.
The fixed functions in the parent attempt remain unchanged.
The thresholds may be chosen to make
the union of the bad roots an arbitrarily small fixed fraction
of $|R|$.  This proves the asserted uniformity and support
convention.

Finally, the estimates have the required relation to the original
cell minimum.  The subadditivity of $x\mapsto\min(\eta,x)$ removes
absolute numerical errors and the critical unpowered terms.  For
the remaining cell sum use
\[
 \min(\eta,x)\le\eta^{1-p}x^p.
\]
The group sum and the prescribed coupled bands are formed before
this inequality is used.  Thus every power cost gains
$\eta^{1-p}$, whereas Equation~\eqref{eq:src-19} is used exactly
as an absolute estimate.  This completes the conditional side
argument.
\end{proof}
